\section{Limitations and Ethical Considerations}
\label{sec:limitations}

\paragraph{Simulated Environment.}
\asb{} uses mock tools rather than live APIs, enabling reproducibility but missing real-world attack surfaces.
On the full 565-case benchmark, BSR drops across all defenses---including D0 baseline (89.5\%$\to$53.1\%)---due to mock-environment limitations on complex multi-tool cases (Appendix~\ref{sec:extended_results}).

\paragraph{Benchmark Scope.}
Our 565 cases cover 6 attack types and 54 injection techniques, but may not represent the full distribution of real-world injections.
Data exfiltration is overrepresented (53.4\%), potentially biasing aggregate metrics.
Cross-benchmark evaluation (\S\ref{sec:cross_benchmark}) revealed that attack success judgments are tightly coupled to each benchmark's evaluation protocol and cannot be transferred via format conversion---standardized cross-benchmark protocols remain an open challenge.

\paragraph{Adaptive Attack Scope.}
Our adaptive evaluation uses a single attacker model (GPT-4o) with 5 iterations; bypass rates measure resistance to feedback-guided optimization rather than absolute robustness (\S\ref{sec:adaptive}).

\paragraph{CIV Limitations.}
All main-table results use CIV~v2 (v1's counterfactual signal was net-harmful).
CIV's provenance signal faces a fundamental tension: legitimate cross-tool data flows and attacker-driven exfiltration can be structurally indistinguishable without task-level intent modeling.

\paragraph{D7 Sanitization Finding.}
Our finding that D7 increases ASR by 4\% should be interpreted as benchmark-specific: our payloads use explicit injection markers that the sanitizer removes, but which also serve as anomaly cues for the model's self-defense.
Marker-free attacks may yield different results.

\paragraph{Ethical Considerations.}
\asb{} is designed to improve agent security.
We release attack payloads for reproducibility using synthetic data with no real user information, and follow responsible disclosure practices.
